\section{Completing edges to chambers}
\label{sec:chambers}
The following construction and theorem are due to D\'iaz Ramos, Molinier and
Viruel~\cite[Theorem 4.1]{DMV}. We retain their construction and prove all four
conditions of Theorem~\ref{UR.Restriction}.

\begin{theorem}[Chamber completion]\label{UR.Chambers}
Let $X$ be a finite properly colored triangle-free graph of minimum degree at
least two, with palette $I$ of size $r$. For every edge $e=\{v,w\}$ and every
$i\in I\setminus\{t(v),t(w)\}$, introduce a vertex $u_{e,i}$ of color $i$.
Make
\[
 B_e=\{v,w\}\cup\{u_{e,i}:i\notin\{t(v),t(w)\}\}
\]
a complete graph, retaining all old vertices and adding no other edges.
The resulting graph $J(X)$ is an incidence geometry on the same palette and
restriction gives
\[
 (C(J(X)),K(J(X)))\cong(C(X),K(X)).
\]
It has $|V(X)|+(r-2)|E(X)|$ vertices and the same connected components as $X$.
\end{theorem}
\begin{proof}
Every $B_e$ has exactly one vertex of each color. A clique containing a new
vertex lies entirely in its $B_e$, since that vertex has no neighbors outside
$B_e$. A clique consisting of old vertices has size at most two by
triangle-freeness. A two-vertex clique supplies its edge $e$; a singleton
extends using an incident edge; the empty clique extends because $X$ has an
edge. Hence every clique is contained in a chamber.

A new vertex has degree $r-1$. At an old vertex $v$, the sets
$B_e\setminus\{v\}$, for edges $e$ through $v$, are pairwise disjoint.
Consequently
\begin{equation}\label{eq:degrees}
 \deg_{J(X)}(v)=\deg_X(v)(r-1).
\end{equation}
The hypotheses imply $r\geq2$ and $\deg_X(v)\geq2$, so the old vertices
are exactly those of degree greater than $r-1$. Every correlation recognizes
them. Their induced graph is $X$, and their old colors are unchanged; the
restriction therefore belongs to $C(X)$.

Suppose a correlation fixes every old vertex. Every old color occurs, so it
fixes every color. If a new vertex has old neighbor set $\{v,w\}$ and color
$i$, it must be $u_{\{v,w\},i}$. Thus all new vertices are fixed and
restriction is injective. If $f\in C(X)$ induces $\sigma_f\in\Sym(I)$,
define
\[
 \widetilde f(v)=f(v),\qquad
 \widetilde f(u_{e,i})=u_{f(e),\sigma_f(i)}.
\]
This is a bijection, with inverse defined from $f^{-1}$. It carries each
$B_e$ onto $B_{f(e)}$, and hence preserves adjacency in both directions.
It fixes all colors precisely when $f$ does. Theorem~\ref{UR.Restriction}
now identifies the pairs. When $r=2$ no vertices are added, and this proof
still applies.

There are $r-2$ new vertices per edge, giving the count. Every new vertex is
attached within the component of its edge, and all old edges are retained.
No components are joined or lost.
\end{proof}

The source's degree argument on physical page 7 uses $r$ in place of $r-1$.
Equation~\eqref{eq:degrees} is the actual degree calculation. The strict
separation needed for its theorem is unaffected. The statement concerns this
construction and these hypotheses; it does not assert that they are logically
necessary for every conceivable pair-preserving completion.
